\documentclass[12pt,a4paper]{article}
\usepackage[utf8]{inputenc}
\usepackage[T2A]{fontenc}
\usepackage[russian]{babel}
\usepackage{amsmath,amssymb}
\usepackage{geometry}
\geometry{margin=2.3cm}
\title{Точный весовой детерминант\\граничного роторного кандидата}
\author{}
\date{6 августа 2026 г.}
\begin{document}
\maketitle

\section{Проверяемый оператор}

Для вещественной плоскости, на которой голономия действует поворотом на
\(q\theta\), дзета-детерминант ковариантного лапласиана на окружности
пропорционален
\begin{equation}
 \lambda_q(\theta)=2-2\cos(q\theta).
\end{equation}

\section{Провал казимирного кандидата}

В представлении \(2V_{1/2}\oplus3V_{3/2}\) положительные веса имеют кратности
\begin{equation}
 q=1:\ 5,
 \qquad
 q=3:\ 3.
\end{equation}
Поэтому точный логарифмический вклад равен
\begin{equation}
 \Gamma_{\det}
 =\frac13\left(5\log\lambda_1+3\log\lambda_3\right),
\end{equation}
а не \(\frac83\log\lambda_1\). Разность имеет вид
\begin{equation}
 \Delta\Gamma
 =\log\frac{\lambda_3}{\lambda_1}
 =2\log|1+2\cos\theta|.
\end{equation}
Она зависит от угла и при целевой точке равна приблизительно \(-0.98749\).
Совпадение среднего казимира не гарантирует совпадения полного детерминанта.

\section{Точный простой весовой рисунок}

Точный логарифмический член дают восемь одинаковых вещественных плоскостей
с \(|q|=1\). Такой рисунок соответствует четырём комплексным фундаментальным
дублетам \(SU(2)\), имеющим вещественную размерность \(16\):
\begin{equation}
 \Gamma_{\det}
 =\frac83\log\lambda_1.
\end{equation}

При том же нормированном следе фиксированный заряд требует средней нормы
\begin{equation}
 q^2_{\mathrm{req}}=6.
\end{equation}
Каноническая осевая норма равна \(1\), а полный удвоенный казимир равен
\(3\). Последний даёт обратный коэффициент \(4\), тогда как требуется \(8\).

\section{Вердикт}

Представление \(2V_{1/2}\oplus3V_{3/2}\) закрыто как точное
детерминантное завершение. Четыре фундаментальных дублета точно
воспроизводят логарифмический член, но оставляют двукратный недостаток
фиксированного зарядового вклада.

Следующая проверка должна либо вывести удвоение зарядового сектора из
действия, либо доказать обязательность единой нормировки и закрыть всю
роторную ветвь.

\end{document}